\section{The True Easter Wonder: Lazarus Bound and Christ Unbound}
\label{sec:lazarus-christ}

To retire the modern folded-napkin story without opening the deeper wonder St. John actually built into his Gospel would leave the empty tomb half-read. The Fourth Gospel contains only two burial scenes, and only two occurrences of the Greek word \latin{soudarion}. St. John placed them nine chapters apart so that every attentive reader who enters the garden sepulchre on Easter morning remembers the cave at Bethany.

\begin{massbox}{John 11:43--44 (Douay--Rheims, Challoner revision)}
\textsuperscript{43}\,When he had said these things, he cried with a loud voice: Lazarus, come forth.
\textsuperscript{44}\,And presently he that had been dead came forth, bound feet and hands with winding bands. And his face was bound about with a napkin. Jesus said to them: Loose him and let him go.
\end{massbox}

At Bethany, the four-days-dead Lazarus hears the commanding cry of the Son of God and emerges from the rock---still wearing his grave. His feet and hands remain tied with narrow winding bands (\latin{keiriais}), and his face remains bound around with a head-cloth (\latin{soudari\={o} periededeto}). Blinded and hobbled in his own burial wrappings, Lazarus cannot even walk home until bystanders step forward at Jesus' command: ``Loose him and let him go.''

The Fathers saw rich meaning in Lazarus's helplessness. St. John Chrysostom (preserved in St. Thomas Aquinas's \emph{Catena Aurea} on John 11:44) observes that Lazarus came forth bound ``that none might suspect that he was a mere phantom,'' and that walking out with bound feet was itself a second miracle alongside his return to life; Jesus ordered the bystanders to loose him so that by touching him with their own hands they would be certain it was the same dead man they had buried four days earlier. St. Augustine (\emph{Tractate 49 on John}, \S 24) reads the scene sacramentally: the sinner buried under the heavy stone of evil habit hears the loud voice of grace and comes forth by confession, yet remains bound until the ministers of the Church exercise the power of the keys---``Whatsoever ye shall loose on earth shall be loosed in heaven'' (Matt 18:18). Cornelius a Lapide adds Augustine's further reflection on the \latin{soudarion} covering Lazarus's eyes: in this mortal life our vision remains veiled, until the veilings are lifted and we see God face to face.

Beneath every moral and sacramental meaning lies a stark theological fact about death itself, stated with doctrinal precision by the Catechism of the Council of Trent on the Fifth Article of the Creed:
\begin{quote}
``These words of the Apostle [\emph{Christ is risen from the dead, the first fruits of them that sleep}, 1 Cor 15:20] are to be understood of a perfect resurrection, by which we are resuscitated to eternal life, being no longer subject to death. In this resurrection Christ the Lord holds the first place; for, if we speak of resurrection, that is, of a return to life, subject to the necessity of again dying, many were thus raised from the dead before Christ; all of whom, however, were restored to life to die again; but Christ the Lord, having conquered death, rose again to die no more, according to this clear testimony of the Apostle: `Christ rising again from the dead, dieth now no more, death shall no longer have dominion over him.'''\endnote{Catechism of the Council of Trent, Part I, Article V, trans. J. Donovan (1829), p. 54, quoting Rom 6:9. Compare Douay--Rheims Rom 6:9 verbatim: ``Knowing that Christ, rising again from the dead, dieth now no more. Death shall no more have dominion over him.''}
\end{quote}

Lazarus came out of his tomb still wearing his \latin{soudarion} because he was returning to mortal life; one day he would lie down again, and friendly hands would bind his face in a shroud once more. Jesus did not return to mortal life. He passed through death into the immortal glory of the Father. No bystander stood at the garden door to untie Him. No human hand had to strip the myrrh-stiffened linens from His limbs or lift the face-cloth from His eyes. When Simon Peter stepped into the rock-cut chamber at dawn, and when the beloved disciple followed him across the threshold, the \latin{othonia} lay quietly on the stone and the \latin{soudarion} that had covered His head lay rolled up apart in one place---not left as a household table code to tell a servant that a meal had paused, but left behind forever because the One who had lain there would never need burial clothes again.

St. Gregory the Great perceived the quiet majesty of that empty stone chamber when he preached on these very verses in Easter Week: the body-cloths belong to our shared mortal frame, yet the head-cloth rolled up apart---whose beginning and end cannot be seen---turns the mind upward to the eternal Godhead of the One who chose death freely and rose by His own divine power. And every year across the Easter Octave, the Church gathers that witness into song in \emph{Victimae paschali laudes}, placing the empty cloths beside the angels as the herald of an unbroken triumph: \latin{Angelicos testes, sudarium et vestes. Surrexit Christus spes mea.} The linen cloths do not whisper ``I stepped away from the table''; they bear silent, orderly witness beside the open rock that death has lost its prisoner forever.
